%===============================================================================
% MINIMAL EXAMPLE: Using Gimmatek LaTeX Template
%===============================================================================
% This is the simplest possible document using the template system.
% Copy this file as a starting point for your own documents.
%===============================================================================

\documentclass[12pt,a4paper]{gimmatek-report-template}

% Include template packages from parent directory structure
% Adjust path if your structure is different
\makeatletter
\def\input@path{{../packages/}}
\makeatother
\RequirePackage{gimmatek-report-frontmatter}

%-------------------------------------------------------------------------------
% Document Metadata
%-------------------------------------------------------------------------------
\doctitle{Minimal Example Document}
\doctitlepage{Minimal Example Document}
\docsubtitle{A Simple Demonstration}
\docversion{V1.0}
\doctype{Example Document}

%-------------------------------------------------------------------------------
% Abstract
%-------------------------------------------------------------------------------
\docabstract{%
This is a minimal example demonstrating the Gimmatek LaTeX Template system.
It shows how to create a professional document with just a few lines of code.

The template handles all front matter generation automatically, including
table of contents, list of tables, and list of figures.
}

%===============================================================================
% DOCUMENT BODY
%===============================================================================
\begin{document}

%-------------------------------------------------------------------------------
% Front Matter - Just 3 commands!
%-------------------------------------------------------------------------------
\frontmatter
\pagenumbering{roman}
\makefrontmatter  % ← This does all the magic!

%-------------------------------------------------------------------------------
% Main Content
%-------------------------------------------------------------------------------
\mainmatter

\chapter{Introduction}

This is the first chapter of the document. Notice how simple the structure is!

\section{Purpose}

The purpose of this example is to demonstrate:

\begin{itemize}
\item Minimal setup required
\item Automatic front matter generation
\item Professional appearance
\item Easy content creation
\end{itemize}

\section{Key Features}

\subsection{Automated Front Matter}

With just one command (\texttt{\textbackslash makefrontmatter}), you get:

\begin{enumerate}
\item Abstract with Chinese title (摘要)
\item Table of Contents (目錄)
\item List of Tables (表目錄)
\item List of Figures (圖目錄)
\item Correct page numbering
\end{enumerate}

\subsection{Professional Styling}

The template provides professional styling out of the box.

\chapter{Using Tables}

\section{Example Table}

Here's a simple table using the booktabs style:

\begin{table}[htbp]
\centering
\begin{tabularx}{\textwidth}{l X}
\toprule
\textbf{Feature} & \textbf{Description} \\
\midrule
Auto TOC & Table of contents generated automatically \\
Auto LOT & List of tables generated automatically \\
Auto LOF & List of figures generated automatically \\
Styling & Professional Stanford-inspired design \\
\bottomrule
\end{tabularx}
\caption{Template Features}
\label{tab:features}
\end{table}

As shown in Table~\ref{tab:features}, the template provides many automated features.

\chapter{Conclusion}

This minimal example demonstrates how easy it is to create professional
documents using the Gimmatek LaTeX Template system.

\section{Next Steps}

To create your own document:

\begin{enumerate}
\item Copy this file as a starting point
\item Update the metadata (title, version, etc.)
\item Write your content
\item Compile with: \texttt{xelatex minimal\_example.tex}
\end{enumerate}

That's it! The template handles the rest.

\end{document}

%===============================================================================
% COMPILATION INSTRUCTIONS
%===============================================================================
%
% To compile this document:
%
% 1. Ensure you have XeLaTeX installed
%
% 2. Set TEXINPUTS to find the template packages:
%    export TEXINPUTS=.:../packages//:
%
% 3. Compile (run twice for references):
%    xelatex minimal_example.tex
%    xelatex minimal_example.tex
%
% 4. View the PDF:
%    open minimal_example.pdf  # macOS
%    xdg-open minimal_example.pdf  # Linux
%
%===============================================================================
